Mit $v = \frac{2\pi r}{T}$ ist die Radialbeschleunigung eines Punktes im Abstand $r$ von der Achse
\begin{align}
 a &= \frac{v^2}{r} = \frac{4\pi^2 r}{T^2}
\end{align}
Bei gleichem Radius ist $a \propto 1/T^2$. Die neue Umlaufzeit ist $T' = \left(1-\pKa\right)T$, also
\begin{align}
 \frac{a'}{a} &= \qKaF = \frac{1}{\left(1-\pKa\right)^2} = \qKa \\
 \frac{a'-a}{a} &= \frac{a'}{a}-1 = \dKa
\end{align}
Die Radialbeschleunigung nimmt um $\result{\dKaP{0}{2}}$ zu.
